% =====================================================================
% 5. Literature and Technology Review (Progress to Date)
% =====================================================================
\section{Literature and Technology Review}
\label{sec:literature}

\subsection{Supervisory control and situation awareness}\label{subsec:supervisory-control-and-situation-awareness}
Sheridan~\cite{sheridan1992telerobotics} defines supervisory control as the monitoring of an automated process with intermittent human intervention, a description that closely matches an engineer observing a long BMS test.
Of the four stages of automation identified by Parasuraman et al.~\cite{parasuraman2000model}, a read-only dashboard performs only information acquisition and analysis, so that all decisions concerning the test remain with the engineer.
Endsley's model of situation awareness~\cite{endsley1995toward} distinguishes between the perception of events, their comprehension and the projection of future states.
These levels correspond to the needs of the different audiences: an engineer is chiefly concerned with perceiving changes in individual cells, whereas a project lead requires comprehension of whether the test is proceeding as planned.

\subsection{Digital twins and test benches}\label{subsec:digital-twins}
Kritzinger et al.~\cite{kritzinger2018digital} define a digital shadow as a representation that receives data automatically from a physical object but does not transmit data back to it.
The dashboard is deliberately designed as a digital shadow, since a system without a write path to the rig cannot disturb a running test.
Research on battery digital twins has concentrated on state estimation and degradation, and a recent review identifies the rapid creation of dashboard interfaces as an unresolved challenge~\cite{naseri2023batterydt}.
Monitoring dashboards with alerting exist for experimental test beds, but they are described rather than evaluated~\cite{darkside2025grafana}.
No study was identified that evaluates monitoring software for verification rigs against principles of alarm management or situation awareness.
\todoin{Run and report a documented Scopus query to support this claim.}

\subsection{Alarm management}\label{subsec:alarm-management}
Wang et al.~\cite{wang2016alarmoverview} report that chattering alarms, which repeatedly activate and clear within a short period, may account for between 10\% and 60\% of all alarm occurrences in process plants.
Since a threshold alert on a noisy cell voltage would behave similarly, the alert engine merges duplicate alerts, applies hysteresis and allows engineers to shelve alerts that have already been acknowledged.
Alert colours follow the aviation convention of red for warnings and amber for cautions~\cite{faa2010alerting}, with blue reserved for advisories.

\subsection{Signal conditioning}\label{subsec:signal-conditioning}
The partner's brief requires that filtering should not misrepresent ``short spikes or noise''.
Moving-average and exponential filters suppress noise at low computational cost but attenuate and delay abrupt changes.
Median and Hampel filters remove isolated outliers while preserving sharp edges, although they may also remove a genuine single-sample spike.
Kalman filters are widely applied to BMS state estimation~\cite{plett2015battery}, but they depend on a system model and careful tuning.
Research on BMS sensor faults has largely evaluated model-based observers using a constant sensor bias~\cite{yu2021sensorfault}; no comparison of simpler filters on BMS data with seeded spikes, drifts and dropouts was found, and RQ2 addresses this gap.

\subsection{Evaluation of dashboards without users}\label{subsec:dashboard-evaluation}
Few~\cite{few2013dashboard} argues that monitoring dashboards should fit on a single screen and avoid decorative elements, which is in tension with the brief's request for a polished customer display.
The Keystroke-Level Model (KLM) predicts expert task completion time from standard operator times, with an error of approximately 21\%~\cite{card1980keystroke}; this accuracy is sufficient for comparing alternative designs.
The feature-congestion measure proposed by Rosenholtz et al.~\cite{rosenholtz2007clutter} estimates visual clutter from a screenshot.
Automated WCAG~2.2~\cite{w3c2023wcag} checks cover at most half of the success criteria~\cite{vigo2013benchmarking}, and a manual checklist is therefore also employed.

\subsection{Existing tools and technology choices}\label{subsec:tech}
Commercial HIL software, such as that supplied by dSPACE and NI, provides instrument panels intended for the engineers conducting a test and is tied to the vendor's hardware.
Grafana is the most credible alternative, as it supports time-series panels and alerting with minimal development effort.
However, it lacks built-in support for shelving and hysteresis, and it is poorly suited to producing distinct views for different audiences; it is therefore adopted as a comparison baseline rather than as the implementation platform.
Table~\ref{tab:tech} summarises the principal technology decisions.

\begin{table}[H]
    \centering\small
    \caption{Technology decisions.}
    \label{tab:tech}
    \begin{tabularx}{\textwidth}{llY}
        \toprule
        Decision & Choice & Rationale \\
        \midrule
        Transport & WebSocket & Lower latency than polling~\cite{pimentel2012websocket} \\
        Rig ingestion & Read-only DB polling & Requires no change on the partner's side \\
        Storage & TimescaleDB (provisional) & SQL support; a single database for runs and samples \\
        Front end & React, canvas charts & Full control over role-specific views \\
        Downsampling & MinMaxLTTB & Preserves peaks at low computational cost~\cite{vandonckt2023minmaxlttb} \\
        \bottomrule
    \end{tabularx}
\end{table}

The remaining review will cover a documented database search for the research gap, the BMS fault types most relevant to the STT, and benchmarks for time-series databases and chart rendering.
